\section{Step 29 Finalize once and preserve the result across close}\label{crc-s29}

At successful EOF, complement the retained raw state, close the descriptor, then place the checksum and success status in their result registers. Keep the checksum in a saved register during close because \texttt{v0} is also the syscall selector. The final branch goes directly to restoration, bypassing temporary cleanup.

\Needspace{16\baselineskip}
\subsection{Write and inspect}
Replace the marked temporary instructions in \texttt{f\_finish}. Retain surrounding labels and markers.

\MARSCodeLabel{MIPS assembly}
\begin{minted}[samepage=true]{mars-idea}
    dli $t0, 0xffffffff
    xor $s5, $s5, $t0
    move $a0, $s4
    li $v0, 16
    syscall
    move $v0, $s5
    move $v1, $zero
    beq $zero, $zero, file_return
    nop
\end{minted}

Use \texttt{s29-1}, \texttt{s29-2}, \texttt{s29-3}. The CLI checks install explicit in-memory files and read/error controls. Ordinary IDE files do not force these short-read or failure conditions.

\subsection{Expected observation}
The complete check file returns unsigned \texttt{cbf43926}/status zero and no open descriptor; the empty file returns zero/status zero.

\Needspace{11\baselineskip}
\subsection{Check your understanding}
Why retain the checksum outside \texttt{v0} while closing?

\begin{itemize}
  \item \textbf{A} The close service needs \texttt{v0} as its selector.
  \item \textbf{B} The polynomial is wider than the register.
  \item \textbf{C} Closing appends a checksum byte.
\end{itemize}
